% X_0^4130(1)* is hyperelliptic: its Jacobian has no rational torsion, so one CM relation D_1 - D_2 between two
% degree-two divisors certified torsion by the dual test is principal, and |D_1| is a g^1_2.
% Tests: tests/Hyperelliptic4130_1.m and tests/_offline/Hyperelliptic4130_1.m.
\documentclass[11pt,reqno]{amsart}
\usepackage[margin=1.15in]{geometry}
\usepackage{amsmath,amssymb,amsthm,mathtools}
\usepackage[colorlinks=true,linkcolor=blue,citecolor=blue,urlcolor=blue]{hyperref}
\theoremstyle{plain}
\newtheorem{theorem}{Theorem}[section]
\theoremstyle{definition}
\newtheorem{remark}[theorem]{Remark}
\DeclareMathOperator{\Jac}{Jac}
\newcommand{\Q}{\mathbb{Q}}
\newcommand{\Z}{\mathbb{Z}}
\newcommand{\F}{\mathbb{F}}
\begin{document}
\title{$X_0^{4130}(1)^*$ is hyperelliptic}
\author{}
\date{\today}
\maketitle

\section{Statement}

Let $Y=X_0^{4130}(1)^*$ be the quotient of $X_0^{4130}(1)$ by its full Atkin--Lehner group $W$, of order $16$.
It has genus three and no further quotient, so it is hyperelliptic if and only if it carries two distinct linearly
equivalent divisors of degree two. Consider
\[
D_1=P_{-35}+P_{-532},\quad D_2=P_{-163}+P_{-595},\quad D_3=P_{-315}+P_{-520},\quad D_4=P_{-340}+P_{-420},
\]
each the sum of two rational CM points of $Y$ (the eight points are rational: their fields of definition have
degree one). They are the fibres of the hyperelliptic map found by the Shimizu screen of the invariant
differentials at the rational CM points; the proof uses them only as divisors.

\begin{theorem}\label{thm:main}
$D_1\sim D_2$. In particular $Y$ is hyperelliptic.
\end{theorem}

\section{Proof}

\emph{$\Jac(Y)(\Q)$ is torsion-free.} The orders $\#\Jac(Y)(\F_p)=L_p(1)$, from the point counts of $Y$ over
$\F_p$, $\F_{p^2}$, $\F_{p^3}$ by the trace formula, have greatest common divisor $1$ over the good primes below
$120$. So a degree-zero divisor on $Y$ is principal as soon as it is torsion.

\emph{$D_1-D_2$ is torsion.} The dual test of \texttt{paper/hyperelliptic-15-38.tex}, Proposition~2.1, at
$M=4\cdot2065=8260$ (\texttt{tests/\_offline/Hyperelliptic4130\_1.m}): a principal part with divisor $D_1-D_2$ in
the dictionary of Guo--Yang pairs to zero with the space $W$ of cusp forms of weight $\tfrac32$ on $\Gamma_0(8260)$
vanishing at every index carrying no CM point; the test checks that $T(9)$ acts on $W$ through a factor of
$x^3+3x^2-3$, the polynomial of the traces of Frobenius of $Y$ at $3$, certifies all six differences $D_i-D_j$,
and refuses $P_{-35}-P_{-532}$ and $P_{-35}+P_{-163}-P_{-532}-P_{-595}$. Two of the points need a word: $P_{-35}$
is a fixed point of $w_{35}$, so it enters the CM cycle $Z(35)$ with weight $\tfrac12$, and $-315=9\cdot(-35)$ is
not fundamental, so $P_{-315}$ shares the cycle $Z(315)$ with $P_{-35}$; the dictionary carries both, and the
linear system of the test separates them because $P_{-35}$ also appears alone in $Z(35)$.

So $D_1-D_2$ is principal; $D_1\neq D_2$ as divisors, so $|D_1|$ is a $g^1_2$ on $Y$. \qed

\begin{remark}
This is the simplest shape of the argument: with no rational torsion, a single certified relation does
everything. Compare \texttt{paper/hyperelliptic-6-649.tex}, where the torsion is bounded by $6$ and the
component group at $3$ is needed to make three fibres suffice.
\end{remark}

\end{document}
